\section{Consequences of model restrictions}
\subsection{Global and local identification}
Assume the model is well specified on the chosen items, with $P_i=q_{\theta^\ast}(\cdot\mid x_i)$ and nonconstant $t=M(\theta^\ast)$.
Proposition~\ref{prop:invariance} gives the exact global optimum set
\begin{equation}\label{eq:thetaset}
 \Theta_\rho=\{\theta\in\Theta:M(\theta)=a\mathbf1+bM(\theta^\ast),\ a\in\R,\ b>0\}.
\end{equation}
Correlation identifies the true predictions exactly when every $\theta\in\Theta_\rho$ satisfies $q_\theta(\cdot\mid x_i)=P_i$ for all $i$.
It identifies the parameter itself only when this set reduces to $\{\theta^\ast\}$, after any stipulated parameter equivalences are removed.
This requires excluding alternative feasible affine copies and alternative shapes with the same means.
Within a well-specified class, CE also identifies the true predictions, by Proposition~\ref{prop:ce}; parameter symmetries can still prevent parameter identification.

Local identification replaces $\Theta$ in this criterion by a neighborhood of $\theta^\ast$.
It does not exclude a distant second optimum.
For an open $d$-dimensional parameter domain and continuously differentiable $M$, the following is a sufficient local condition:
\begin{equation}\label{eq:tangent}
 DM(\theta^\ast)v\in\operatorname{span}\{\mathbf1,t\}\quad\Longrightarrow\quad v=0.
\end{equation}
To see this, define the standardized mean map
\[
 B(\theta)=\frac{M(\theta)-\overline{M(\theta)}\mathbf1}{\|M(\theta)-\overline{M(\theta)}\mathbf1\|}.
\]
For any perturbation $h$ of the mean vector, its derivative under centering and normalization vanishes exactly when $h\in\operatorname{span}\{\mathbf1,t\}$ at $\theta^\ast$.
Thus Eq.~\eqref{eq:tangent} makes $DB(\theta^\ast)$ have full column rank.
Selecting $d$ independent output coordinates and applying the inverse function theorem shows that $B$ is injective in a neighborhood.
Since correlation one is equivalent to $B(\theta)=B(\theta^\ast)$, the local optimum is unique there.
Failure of this first-order condition alone does not prove nonidentification; a singular derivative can coexist with local injectivity.

A family name, including logistic-linear, does not by itself establish Eq.~\eqref{eq:thetaset}'s uniqueness on a finite design.
The number and placement of the items and any shared-parameter symmetries matter.
A claim about identities of functions on an interval is stronger than a claim about equality on finitely many items.

Restricted models can also contain approximate affine copies.
For softmax logits $c v_i(r)$ with fixed $v_i$ and $c\to0$, direct Taylor expansion gives
\begin{equation}\label{eq:weaklogit}
 m_i(c)=\frac{K-1}{2}+\frac cK\sum_r\left(r-\frac{K-1}{2}\right)v_i(r)+O(c^2).
\end{equation}
Indeed $q_i(r;c)=K^{-1}+cK^{-1}(v_i(r)-\bar v_i)+O(c^2)$, and multiplication by $r$ and summation gives Eq.~\eqref{eq:weaklogit}.
On a fixed finite design the remainder is uniform across items.
If the displayed linear coefficient is nonconstant, then for any fixed $a>0$, $\rho(m(c),m(ac))\to1$ as $c\to0^+$.
This establishes approximate scale ambiguity near uniform logits, not an exact additional symmetry of every softmax model.
It supplies no empirical explanation for a fitted model without checking that model's regime.

\subsection{Linear mean models}
\begin{proposition}[Best linear direction]\label{prop:linear}
Let $\phi$ be a random feature vector with positive definite covariance $\Sigma$, let $t$ be a target mean with $\Var(t)=\sigma_t^2>0$, and put $c=\Cov(\phi,t)\ne0$.
For unconstrained linear mean predictions $m=a+w^\top\phi$,
\[
 \rho(w)=\frac{w^\top c}{\sqrt{w^\top\Sigma w}\,\sigma_t}.
\]
Its maximum is attained exactly at $w=\beta\Sigma^{-1}c$, $\beta>0$, with
\[
 \rho_{\max}^2=\frac{c^\top\Sigma^{-1}c}{\sigma_t^2}.
\]
The offset $a$ and positive scale $\beta$ are undetermined.
\end{proposition}
\begin{proof}
Centering removes $a$, and the covariance and variance of $w^\top\phi$ give the first formula.
Cauchy--Schwarz applied to $\Sigma^{1/2}w$ and $\Sigma^{-1/2}c$ gives
\[
 (w^\top c)^2\le(w^\top\Sigma w)(c^\top\Sigma^{-1}c).
\]
Equality requires proportionality of these vectors, and a positive correlation maximum requires a positive proportionality factor.
This yields $w=\beta\Sigma^{-1}c$.
The same vector solves the normal equations for least squares with an intercept.
Its explained variance is $c^\top\Sigma^{-1}c$, giving the stated squared correlation.
\end{proof}

The result also applies to empirical covariances on a fixed item set.
If $c=0$, every nonconstant feasible linear forecast has correlation zero, so there is no distinguished direction.
If forecasts must lie in $[0,K-1]$, the optimum above requires a feasible offset and scale; the unconstrained formula alone does not guarantee feasibility on unbounded feature supports.
Correlation can equal one in this linear class only when the target is an affine function of the features almost surely.
A linear mean head itself provides no distributional shape output.

For comparison, a softmax-linear distribution model $q_\theta(r\mid x)=\softmax(W\phi(x))_r$ has expected CE gradient
\[
 \frac{\partial L}{\partial W_r}=\E\big[\phi(x)(q_\theta(r\mid x)-P(r\mid x))\big].
\]
This follows by applying Eq.~\eqref{eq:gradients} and the chain rule.
A finite unpenalized stationary point matches these feature-weighted category frequencies.
Under misspecification this need not recover the true conditional distribution or each item's mean.
The moment-matching guarantee in Proposition~\ref{prop:exponential} assumes a free scalar natural parameter for the distribution being fitted; arbitrary shared-parameter classes do not inherit that guarantee.

\subsection{Nonlinear mean models}
A nonlinear mean model can represent a nonlinear target $t(x)$ and attain population correlation one.
Its input-to-rating map need not be linear.
If an additional output changes the distribution along a feasible path while preserving all means, the objective is constant on that path.
If shared parameters force every shape change to alter means, correlation can influence shape indirectly, as Proposition~\ref{prop:gradient} states.
Thus the unidentifiable set must be intersected with the feasible model class rather than inferred from a parameter's informal description as a spread parameter.
